\chapter{Operating}

\label{ch:operate}

This is the chapter nobody reads until something has gone wrong, and it is written for that
moment: short procedures, the exact command for each, the reading that tells you whether
it worked, and the failure modes that have actually happened on this product rather than
the ones a design review imagined. The order is the order of the questions an operator
asks --- can I get the data back, what is the store doing to the disk, what changed when I
upgraded, where is the evidence, and why is the daemon not answering.

\section{Backup and restore}

A backup is offline by construction. \texttt{pii backup} opens each store with redb's
exclusive lock and copies the file while holding that lock, so a running daemon makes the
backup \emph{refuse} --- with the path and the words \enquote{stop the daemon first} ---
instead of copying a file with a writer mid-commit. Each copy is then re-opened as a redb
database and hashed, so a torn or truncated copy is an error rather than a silent bad
backup.

\begin{lstlisting}[language=bash]
systemctl --user stop pii
pii backup /backups/pii-$(date +%F) --dir ~/.local/state/pii
systemctl --user start pii
\end{lstlisting}

The command prints one line per file, with these field names and no others:
\texttt{backed up store.redb (\textless{}bytes-> bytes,
format\_version=\textless{}n->, blake3=\textless{}hash\textgreater{})}.
The same facts go into \texttt{MANIFEST.json} inside the destination, and the manifest is
what a restore checks against. The destination must be new or empty.

\begin{center}
\resizebox{\linewidth}{!}{
\begin{tikzpicture}[node distance=9mm and 13mm]
  \node[s] (live) {live stores\\\scriptsize six redb files};
  \node[b, right=of live] (lock) {\texttt{pii backup}\\\scriptsize takes each lock};
  \node[b, right=of lock] (copy) {copy\\\scriptsize through a temp name};
  \node[b, below=of copy] (hash) {re-open + blake3\\\scriptsize a torn copy is an error};
  \node[b, left=of hash] (manifest) {\texttt{MANIFEST.json}\\\scriptsize size, version, hash};
  \draw[f] (live) -- (lock);
  \draw[f] (lock) -- (copy);
  \draw[f] (copy) -- (hash);
  \draw[f] (hash) -- (manifest);
  \draw[bk] (lock) to[bend right=32] node[p,above=2pt,midway]{holds it} (live);
\end{tikzpicture}}
\captionof{figure}{A backup that cannot tear. The red edge is the normal case on a live host:
the open fails, the command refuses, and no file is written.}
\end{center}

Two details decide whether the backup is worth having.

First, the backup set is six files: \texttt{store.redb}, \texttt{memory.redb},
\texttt{approval.redb}, \texttt{heal.redb}, \texttt{channel.redb} and
\texttt{timer.redb}. \texttt{pii health} reports seven stores, and the seventh ---
\texttt{queue.redb}, where the work rows live --- is not in the backup set at this
revision. If you rely on the queue, copy that file with the daemon stopped as well, or
record the rows somewhere you can rebuild them from.

Second, \texttt{--dir} names the directory that holds the stores, and the durable pair
(channel, timer) defaults to \texttt{\textless{}state dir\textgreater{}/durable}, not to the
directory beside the
main store. A backup run with \texttt{--dir \textless{}state dir\textgreater{}} therefore
captures the main store,
memory, approvals and heal, and reports nothing for channel and timer because those files
are one level down. Run a second backup with \texttt{--dir \textless{}state
dir\textgreater{}/durable} if cursors
and timers matter to you --- and they matter the first time a channel replays from the
start after a restore.

Restore verifies every file's hash against the manifest before touching the target, never
deletes, and refuses a store that is in use.

\begin{lstlisting}[language=bash]
systemctl --user stop pii
pii restore /backups/pii-2026-11-04 --dir ~/.local/state/pii           # refuses if a target exists
pii restore /backups/pii-2026-11-04 --dir ~/.local/state/pii --force   # moves the old store aside
systemctl --user start pii
\end{lstlisting}

\texttt{--force} does not delete: it moves the existing store to
\texttt{\textless{}name\textgreater{}.redb.old-\textless{}ms\textgreater{}} and leaves the
decision about the old file to you. That is
deliberate --- the one file you are most likely to want after a bad restore is the store
the restore replaced.

\section{Retention}

Retention is the only code in the product that deletes run data, so its rules are worth
knowing before you need them. The policy has three numbers and three guards.

\begin{center}
\begin{tabular}{@{}p{0.26\textwidth}p{0.62\textwidth}@{}}
	\toprule
	\textbf{Setting} & \textbf{Default and meaning} \\
	\midrule
	run age & 30 days: a finished run's header, journal and snapshots go past this age \\
	node-event age & 7 days: a finished run's node-call events go past this age \\
	rows per sweep & 1\,000: the sweep's whole budget, so a first run on a large store
	cannot hold a write transaction for minutes \\
	\bottomrule
\end{tabular}
\end{center}

The guards are the part that matters:

\begin{enumerate}
	\item A run whose status is \texttt{Running} or \texttt{Suspended} is never removed,
	however old it is, because the daemon may still resume it.
	\item A row that cannot be dated is kept. A run created before the clock was recorded
	carries \texttt{created\_ms} of zero, and deleting on a missing timestamp would destroy
	exactly the runs an operator is looking for.
	\item The sweep is bounded per call, so it is a trim rather than a migration.
\end{enumerate}

One sweep runs when the daemon starts, and \texttt{pii store prune} runs one on demand. The
command prints \texttt{pruned N rows} and exits 0; the store also records a
\texttt{store.retention.swept} event carrying the count, but only when something was
removed. A sweep that removes nothing is silent, which is why the count on stdout is the
number to read.

\section{What retention actually removes}

The three rules have an order of precedence, and the order is not the one people assume.
Age selects candidates; a live run is excluded; an undatable row is excluded; then the
per-sweep budget truncates whatever remains.

\begin{center}
\resizebox{\linewidth}{!}{
\begin{tikzpicture}[node distance=8mm and 12mm]
  \draw[->, thick, gray] (0,0) -- (13,0) node[right, font=\scriptsize] {age};
  \fill[orange!25] (0,0.35) rectangle (2.2,0.75);
  \node[font=\scriptsize, above=2pt] at (1.1,0.78) {node events: 7 days};
  \fill[blue!20] (0,0.9) rectangle (6.4,1.3);
  \node[font=\scriptsize, above=2pt] at (3.2,1.33) {runs: 30 days};
  \draw[dashed, red!60!black, thick] (6.4,0) -- (6.4,1.5);
  \node[font=\scriptsize, red!60!black, below=2pt] at (6.4,-0.05) {sweep boundary};
  \node[b, fill=green!8, anchor=west] (live) at (7.2,1.35)
    {Running / Suspended\\\scriptsize kept however old};
  \node[b, fill=gray!8, below=of live] (undated) {no clock\\\scriptsize kept and reported};
  \node[b, below=of undated] (budget) {budget\\\scriptsize 1\,000 rows per sweep};
  \draw[bk] (6.4,0.7) -- (live.west);
\end{tikzpicture}}
\captionof{figure}{Three rules, in order of precedence: age bounds the sweep, a live run is
never swept, and a row that cannot be dated is kept. The budget is last, so a large store trims
over several sweeps rather than one long transaction.}
\end{center}

A run that will never finish is therefore never collected. If you see one, park or cancel
it deliberately --- \texttt{pii run abandon} with a reason, or \texttt{pii run cancel} ---
and retention will collect it thirty days after it ended. Retention will not make that
decision for you, and should not: the difference between trimming and losing work in flight
is this guard.

\section{Compaction is not retention}

These two get confused, and the confusion costs a maintenance window that reclaims nothing.
Compaction reclaims free pages \emph{inside} a file; retention removes rows. redb never
returns space to the filesystem except at the end of the file, and only compaction moves
that boundary.

The measured pairs from this product's own runbook are the honest expectation:

\begin{center}
\begin{tabular}{@{}p{0.50\textwidth}p{0.38\textwidth}@{}}
	\toprule
	\textbf{Store state} & \textbf{Before $\rightarrow$ after} \\
	\midrule
	fragmented: 20 sessions of 400\,KB, every other one deleted &
	13\,111\,296 $\rightarrow$ 6\,623\,232 bytes \\
	a store after 400 runs through a daemon & 1\,048\,576 $\rightarrow$ 786\,432 bytes \\
	idle auxiliary stores (memory, approval, heal, channel, timer) &
	1\,056\,768 $\rightarrow$ 36\,864 bytes each \\
	\bottomrule
\end{tabular}
\end{center}

So on a store that has only grown with runs, expect the tail slack back and nothing more.
The 400-run row is the one to remember: it is a two-hundred-kilobyte improvement on a
megabyte file, which is not a reason to schedule an outage. The auxiliary row is the reason
to run compaction at all --- five stores that were preallocated and never wrote shrank by
about 97\,\%.

\section{The restore drill}

A backup nobody has restored is a hope. The drill is short, it uses a scratch state
directory, and it is the only way to know that the manifest, the hashes and the binary's
format version all agree.

\begin{center}
\resizebox{\linewidth}{!}{
\begin{tikzpicture}[node distance=8mm and 12mm]
  \node[s] (store) {live store};
  \node[b, right=of store] (stop) {stop the daemon};
  \node[b, right=of stop] (copy) {\texttt{pii backup}};
  \node[b, below=of copy] (scratch) {scratch state dir\\\scriptsize \texttt{PII\_STATE\_DIR}};
  \node[b, left=of scratch] (probe) {scratch daemon\\\scriptsize health, run ls, trace};
  \node[v, below=of probe] (verdict) {restorable?};
  \draw[f] (store) -- (stop);
  \draw[f] (stop) -- (copy);
  \draw[f] (copy) -- (scratch);
  \draw[f] (scratch) -- (probe);
  \draw[f] (probe) -- (verdict);
  \draw[bk] (verdict.west) to[bend left=22] node[p,left=2pt,midway]{no} (copy.south west);
\end{tikzpicture}}
\captionof{figure}{A backup is a backup once it has been restored into a scratch state directory
and read there. Restoring over the live store destroys the evidence you would need if the
restore is wrong.}
\end{center}

\begin{lstlisting}[language=bash]
mkdir -p /projects/_tmp/restore-drill
pii restore /backups/pii-2026-11-04 --dir /projects/_tmp/restore-drill
PII_STATE_DIR=/projects/_tmp/restore-drill pii store stats
PII_STATE_DIR=/projects/_tmp/restore-drill pii graph trace <run_id> \
  --db /projects/_tmp/restore-drill/store.redb --terminal
\end{lstlisting}

Three readings close the drill. \texttt{store stats} prints a row count per table, so an
empty journal is visible immediately. A trace on a known run id prints that run's terminal
value, which proves the journal survived the round trip. And a daemon started on the
scratch directory with a different runtime dir proves the binary can open the restored
format version --- the one thing a restore cannot repair is a store written by a newer
build.

\section{Upgrading}

An upgrade is four steps and one expected non-zero exit code. The exit code is the part
that alarms people: immediately after installing a new binary and before restarting the
service, \texttt{pii health} exits 3 and prints a \texttt{DRIFT} line, because the client
is the new build and the daemon is still the old process. That is the instrument working.
It returns 0 only after the restart.

\begin{lstlisting}[language=bash]
systemctl --user stop pii
pii backup /backups/pii-pre-upgrade --dir ~/.local/state/pii
bash deploy/install-pii.sh                                      # card-gated, moves the old aside
pii health                                                      # expect 3 + DRIFT
systemctl --user restart pii
pii health && pii --version                                     # expect 0, same revision
\end{lstlisting}

\begin{center}
\resizebox{\linewidth}{!}{
\begin{tikzpicture}[node distance=8mm and 12mm]
  \node[b] (stop) {stop the daemon};
  \node[b, right=of stop] (backup) {backup};
  \node[b, right=of backup] (install) {install};
  \node[v, right=of install] (drift) {health?};
  \node[b, below left=of drift] (three) {3: expected\\\scriptsize new client, old daemon};
  \node[b, below right=of drift] (zero) {0: already the same build};
  \node[b, below=of three] (restart) {restart the service};
  \node[b, below=of restart] (verify) {\texttt{pii health} = 0\\\scriptsize both revisions equal};
  \draw[f] (stop) -- (backup);
  \draw[f] (backup) -- (install);
  \draw[f] (install) -- (drift);
  \draw[f] (drift) -- node[p,left=2pt,midway]{usual} (three);
  \draw[f] (drift) -- node[p,right=2pt,midway]{rare} (zero);
  \draw[f] (three) -- (restart);
  \draw[f] (restart) -- (verify);
\end{tikzpicture}}
\captionof{figure}{The upgrade sequence, with the reading that is correct and looks wrong. A
\texttt{3} here is not a failed upgrade; it is the proof that the install replaced the binary
and the service has not yet been restarted.}
\end{center}

The state directory is forward-compatible only in one direction. Every store carries a
format version stamped on create or first open, and a binary that meets a newer version
refuses to open it:

\begin{lstlisting}[language=bash]
error: invalid: store '/var/lib/pii/memory.redb' has format version 99, which is newer
than the version 1 this binary understands; upgrade pii or use a different store path
\end{lstlisting}

The daemon exits 2 and writes nothing to the file. The remedy is a binary that understands
that version, or a restore of a backup taken by the older build --- never an edit of the
stamp.

Rollback is the install script's inverse and it moves rather than copies:
\texttt{bash deploy/rollback-pii.sh} stops the unit, restores the newest
\texttt{\$PREFIX.old-\textless{}epoch\textgreater{}} contents, keeps what it replaced under
\texttt{\$PREFIX.rolledback-\textless{}epoch\textgreater{}}, and prints each file. Then
\texttt{daemon-reload} and \texttt{restart}, then \texttt{pii health} for a 0.

\subsection*{The gate the card quotes}

No install runs without a deploy card for the revision being installed. The card is a JSON
file with five required fields --- \texttt{incident} (id, fingerprint, count),
\texttt{change} (branch, revision, diff stat), \texttt{proof} (red on base, green on
branch), \texttt{gate} (verdict) and \texttt{rollback} (rehearsal) --- and a card missing
any of them is refused with the field named and no approval path printed. Approval is bound
to the revision the card names, and every decision appends one line to the operator record
(\texttt{PII\_ACTS}, default \texttt{\_tmp/autonomy/operator-acts.log}).

The gate verdict has a condition attached that is easy to miss: a deploy run must quote a
\emph{fresh-tree} gate, and a \texttt{mode=routine} verdict is not release evidence. A
routine gate runs in the working tree against a warm target; a fresh gate exports
\texttt{HEAD} and builds in a throwaway directory. When the two disagree, believe the fresh
one --- a gate can be green in a tree and red in a clean clone.

\section{Logs, audit, and the death that leaves no line}

The daemon writes events as JSON lines, one per event, each stamped with \texttt{ts\_ms}.
The log is off unless a path is set. Three keys control it, and each has a flag that
overrides it.

\begin{longtable}{@{}p{0.24\textwidth}p{0.30\textwidth}p{0.38\textwidth}@{}}
	\toprule
	\textbf{Key} & \textbf{Default} & \textbf{Flag} \\
	\midrule
	\endhead
	\texttt{[core].log\_file} & absent: no log file & \texttt{--log-file} \\
	\texttt{[core].log\_rotate\_bytes} & 10\,485\,760 (10 MiB); 0 never rotates &
	\texttt{--log-rotate-bytes} \\
	\texttt{[core].log\_keep\_files} & 5 rotated files besides the live one &
	\texttt{--log-keep-files} \\
	\bottomrule
\end{longtable}

Rotation renames rather than truncates: \texttt{pii.log} becomes \texttt{pii.log.1},
\texttt{pii.log.1} becomes \texttt{pii.log.2}, and the oldest is discarded, so disk use is
bounded by about \texttt{(keep + 1) * rotate\_bytes}. The file is rotated \emph{before} a
line would push it past the limit, so a single pathological line can exceed the bound once.

Denials are separate and always on: they go to \texttt{audit.log} in the store directory,
rotated the same way, whatever the log file is set to. If you are investigating a
permission problem and the daemon log is clean, read \texttt{audit.log}.

The one reading that matters is the last line. A log that ends on \texttt{daemon.started}
with no \texttt{daemon.stopped} means the process died hard --- a \texttt{kill -9}, an OOM
kill, a machine reset. \texttt{pii health} agrees: \texttt{since\_deploy.starts} and
\texttt{since\_deploy.hard\_deaths} are counters kept across restarts, so a host that has
restarted eleven times and died hard twice says so in two lines.

\section{Disk is owned, bounded and visible}

Disk needed its own command because the obvious answers are about the wrong filesystem. The
measured incident: the container's \texttt{/} is a 288\,GB overlay that reached 16\,GB
free (95\,\%), while \texttt{/projects} --- a separate 594\,GB volume --- had 147\,GB free.
The gate refuses to run below a 12\,GB floor on \texttt{/}, so work stopped, and
\texttt{cargo clean} on the repository freed 152\,GB on the volume that was never short.

\begin{center}
\resizebox{\linewidth}{!}{
\begin{tikzpicture}[node distance=10mm and 16mm]
  \node[s] (overlay) {\texttt{/}\\\scriptsize 288 GB overlay};
  \node[b, right=of overlay] (floor) {12 GB floor\\\scriptsize measured here};
  \node[b, right=of floor] (refuse) {gate refuses};
  \node[s, below=of overlay] (vol) {\texttt{/projects}\\\scriptsize 594 GB volume};
  \node[b, right=of vol] (artefacts) {artefacts\\\scriptsize cache, targets, trash};
  \node[b, right=of artefacts] (lie) {the floor was elsewhere};
  \draw[f] (overlay) -- (floor);
  \draw[f] (vol) -- (artefacts);
  \draw[f] (floor) -- (refuse);
  \draw[bk] (artefacts) -- node[p,above=2pt,midway]{filled this} (lie);
\end{tikzpicture}}
\captionof{figure}{The mismatch that hid the growth: the floor was measured on \texttt{/}, the
artefacts landed on \texttt{/projects}. Freeing the volume looked like progress and was worth
nothing.}
\end{center}

\texttt{pii disk} prints the artefact table and \emph{both} free-space lines, which is the
whole fix for that class of confusion. Four mechanism families fill the disk, and each has
an owner now:

\begin{center}
\begin{tabular}{@{}p{0.30\textwidth}p{0.62\textwidth}@{}}
	\toprule
	\textbf{Mechanism} & \textbf{Owner and bound} \\
	\midrule
	the warm gate cache & ceiling of 40\,GiB
	(\texttt{FRESH\_GATE\_CACHE\_MAX\_GB}); trimmed before a build, never while a build is
	live. It once reached 105.8\,GiB because nothing bounded it. \\
        one-off experiment targets & \texttt{\_tmp/targets/\textless{}purpose\textgreater{}}
with a TTL of 7 days
	(\texttt{DISK\_TTL\_DAYS}), removed by \texttt{cargo xtask disk --gc} \\
	trash & dated directories with a retention window, not a build artefact \\
	failed gate exports & removed after \texttt{DISK\_EXPORT\_TTL\_HOURS} unless a gate
	still owns one \\
	\bottomrule
\end{tabular}
\end{center}

\texttt{cargo xtask disk --gc} refuses to remove anything while a cargo or rustc is live
and prints why, which is the guard that keeps the cleanup from becoming the incident. The
12\,GB floor is not lowered to make a gate pass; if the floor is what stops you, the answer
is space.

\section{When the daemon is wedged}

The sequence is short and it is ordered so that no step destroys evidence the next one
needs.

\begin{enumerate}
	\item \texttt{pii health --json} with a short timeout. No answer within a few seconds
	means the daemon is wedged or dead, and exit 2 means nothing is listening at all.
	\item Read the last lines of the log file before stopping anything. A
	\texttt{daemon.started} with no \texttt{daemon.stopped} is a hard death; the last
	non-start event names what it was doing.
	\item \texttt{systemctl --user stop pii}. Shutdown waits at most five seconds for
	in-flight blocking work, then abandons it and exits anyway. If something was stuck,
	stderr carries exactly one line:
	\enquote{pii: shutdown drain exceeded 5s; abandoning in-flight blocking work}. That
	line is stderr only --- it is in the journal, not in the log file.
	\item If systemd had to SIGKILL it, nothing is lost structurally: the store is a
	transactional redb file, and the daemon reconciles interrupted runs on boot. Start it
	and read \texttt{pii health}.
	\item A run that was mid-flight comes back as \texttt{Running} or \texttt{Failed}, with
	a \texttt{NodeInterrupted} event in its journal for the node that was left in flight.
	Read the trace before resuming it.
\end{enumerate}

\section{Failure modes that have actually happened}

\begin{longtable}{@{}p{0.28\textwidth}p{0.64\textwidth}@{}}
	\toprule
	\textbf{Symptom} & \textbf{Cause and what to check} \\
	\midrule
	\endhead
	Every graph run fails in a second, with no output &
	A model id the provider does not serve. The old id had not existed for months and the
	graph died at spawn; run \texttt{pii graph lint} and read the model in the graph. \\
	The daemon answers, but the store is empty & Two different \texttt{PII\_STATE\_DIR}
	values --- one for the daemon, another for your shell --- or a daemon whose
	\texttt{--store} points at the runtime directory. Read the path in \texttt{pii health}. \\
	A row stays \texttt{Running} for ever & A callee refused and no edge carried the
	refusal. At this revision the queue wires \texttt{refused} and \texttt{unverified} by
	name; an older graph dead-ends inside the splice. Trace the run; the refusal is in the
	journal. \\
	The coder stops and says no write tool is available & The graph's \texttt{caps} omit
	\texttt{fs\_write}, so the node may not hold it. That is the capability ladder working,
	not a bug. \\
	The gate refuses with a disk message on a full disk & The floor now names the filesystem
	it measured. A message that names only \texttt{/} while the artefacts land elsewhere is
	an older binary. \\
	\texttt{backup} refuses with \enquote{cannot be opened exclusively} & The daemon is
	running and holds the lock. Stop it; the refusal is the no-tear guarantee, not a fault. \\
	\texttt{health} prints \texttt{UNVERIFIED: daemon build revision is unknown} & The
	running binary was built without its revision --- an export build that did not receive
	\texttt{PII\_BUILD\_REVISION}. Rebuild with the installer. \\
	A channel replays its backlog after a restore & The channel cursor lives in
	\texttt{durable/channel.redb}, and a backup of the store directory does not include it.
	Back that directory up too. \\
	\bottomrule
\end{longtable}

\section{The first five minutes of an incident}

Order matters more than completeness. Five checks establish whether you are looking at a
process problem, a data problem or an environment problem, and each takes seconds.

\begin{longtable}{@{}p{0.06\textwidth}p{0.30\textwidth}p{0.56\textwidth}@{}}
	\toprule
	\textbf{Step} & \textbf{Command} & \textbf{What it settles} \\
	\midrule
	\endhead
	1 & \texttt{pii health} & Is anything answering, and is it this build? Exit 2 is a
	process problem; exit 3 is a deployment problem. \\
	2 & \texttt{pii run ls --status running} & Is work in flight? A run that should have
	finished is a graph problem, not a daemon problem. \\
	3 & \texttt{pii disk} & Is the host out of space on the filesystem the artefacts use?
	A gate or a write failing at 95\,\% looks like a code defect. \\
	4 & the log's last ten lines & Did it die hard, and what was the last thing it did? \\
	5 & \texttt{pii store stats} & Is the store growing at a rate that explains the
	symptom, and does the row count match what you expect? \\
	\bottomrule
\end{longtable}

The order is not arbitrary. Checking the store before the disk gives you a row count you
cannot interpret; checking runs before health gives you a list from a daemon that may not
be the binary you think. Health first, because every later reading depends on knowing which
process produced it.

\section{Two rules that keep a host boring}

The rest of this manual is detail; these two rules prevent most of the incidents in it.

\textbf{Never point a second process at a live state directory.} The stores are exclusive,
the daemon claims the durable pair before it binds, and a second daemon will be refused ---
but a CLI command with a different \texttt{PII\_STATE\_DIR} will not be refused, it will
quietly read a different store and report an empty world. If a reading surprises you, the
first question is which store it came from.

\textbf{Never delete what you can move.} The installer moves old binaries aside, the
restore moves the replaced store aside, retention keeps undatable rows, and compaction
refuses while a writer holds the file. Every one of those behaviours exists because the
alternative destroyed something nobody could reconstruct. When a procedure in this chapter
seems to be missing a cleanup step, that is the step it is missing on purpose.
